\SourceLeaf{241}%
% A544-P0241-D01
aliquam $x^\nu-1$ dividunt, existente $\nu'<\nu$, iam inventos esse, reliquosque investigare proponi. Hi autem omnes in hac expressione comprehendi possunt $(P,\rho^k)$, si $P$ sit unus ex ipsis et pro $k$ omnes numeri minores quam $\nu$ ad ipsumque primi substituantur.

% A544-P0241-D02
In Cap. \textsc{vi} ostendimus, quomodo radices primae aequationis $x^\nu=1$ ita in classes discerpi possint, ut, omnibus per alicuius potestates expressis, eadem in classes distributio habeatur, quaecunque radix prima pro hac basi accipiatur; \emph{periodos} huiusmodi radicum complexus vocavimus. Iam patet, functiones $x,x^\alpha,x^\beta,x^\gamma$ etc., designantibus $\alpha,\beta,\gamma$ etc. omnes numeros ad $\nu$ primos, simili modo in periodos resolvi posse, quamque periodum maiorem rursus in minores donec tandem ad periodos formae $x^k,x^{kp},x^{kpp}\dotstwo x^{kp^{m-1}}$ perveniatur. Hoc ita facto patet

% A544-P0241-D03
$1^{\mathrm{o}}$ Quoniam periodus quaeque ex huiusmodi periodis minimis $x^k+x^{kp}+\text{etc.}$ composita est, si per quamcunque functionem primam $m$ dimensionum dividatur, residuum fore numerum.

% A544-P0241-D04
$2^{\mathrm{o}}$. Quum omnes periodi termini semper ad hanc formam reduci queant $x^{\kappa\pdot a^\alpha b^\beta c^\gamma}\dotstwo$, ubi $\kappa.\ a,b,c\dotstwo$ sunt numeri determinati, pro $\alpha,\beta,\gamma\dotstwo$ autem omnes valores substitui possunt; patet, periodum in se ipsam mutari, si pro $x$ substituatur $x^k$ et $k$ sit formae $a^\alpha b^\beta c^\gamma\dotstwo(\mathrm{mod.}\,\nu)$, unde facile perspicitur omnes functiones $P,(P\;\rho^k)$ etc., designante $k$ huiusmodi numerum, si periodus per eas dividatur, idem residuum dare.

% A544-P0241-D05
$3^{\mathrm{o}}$. Quare periodus subducto tali residuo per productum ex omnibus functionibus $(P,\rho^k)$ dividi poterit.

% A544-P0241-D06
\Article{363}
Summa rei in hoc vertitur, ut haec residua determinentur. Primo quaeratur residuum, quod periodus maxima per productum ex omnibus functionibus primis idoneis dabit. Si hoc productum sit
\[
\equiv x^\lambda-Ax^{\lambda-1}+\text{etc.}
\]
erit residuum hoc $\equiv A$. Huius autem producti forma facile invenitur et ex Cap. \textsc{vi} sequitur esse $A=0$, si $\nu$ per quadratum dividi possit, contra esse $A$ aut $=+1$ aut $=-1$, prout multitudo factorum primorum numeri $\nu$ sit par aut impar.

% A544-P0241-D07
Iam resolvatur haec periodus maxima in periodos inferiores repraesententurque periodi cuiusvis termini per $x^{kp^\pi u}$, ita ut $k$ in quavis periodo sit numerus
\SourceLeaf{242}%
% A544-P0242-D01
determinatus, pro diversis vero variabilis, $\pi$ et $u$ autem in quavis periodo variabiles, eos autem valores, quos in aliqua periodo habent, etiam in reliquis adipisci possint. Supponatur aliquantisper aliqua functio prima $P$ pro basi sitque residuum, quod periodi $\Sigma x^{p^\pi u}$, $\Sigma x^{k'p^\pi u}$ etc. per eam divisae praebent respective $A,A'$ etc., erit $\Sigma x^{p^\pi u}-A$ per productum ex omnibus functionibus $(P,\rho^u)$ divisibilis. $\Sigma x^{k'p^\pi u}-A'$ per productum ex omnibus functionibus $(P,\rho^{k'u})$ etc. etc. At facile liquet, quantitates $A,A'$ etc. esse radices congruentiae datae. Scilicet sint periodi radicum aequationis $x^\nu=1$ periodis praecedentibus correspondentes radices aequationis $Q=0$, erunt $A,A'$ etc. radices congruentiae $Q\equiv0$. Namque erit
\[
\begin{aligned}
A+A'+\text{etc.}&\equiv\text{summae periodorum},\\
AA+A'A'+\text{etc.}&\equiv\text{summae quadratorum periodorum}
\end{aligned}
\]
etc. etc. Calculus enim prorsus similis erit ei, quem Cap. \textsc{vi} exposuimus, si pro $\rho$ substituatur $x$, quoniam etiam hic poni potest pro $x^\nu$ unitas, uti illic pro $\rho^\nu$.

% A544-P0242-D02
Inventis radicibus $A,A'$ etc. aliqua pro residuo periodi $\Sigma x^{p^\pi u}$ eligatur et inde reliquarum residua simili modo uti Cap. \textsc{vi} ordinentur. Namque illud etiam hic arbitrio relinquitur, quum functio $P$ sit prorsus hactenus indeterminata. Calculus sequens omnino analogus est ei, quem Cap. \textsc{vi} pertractavimus, singula exponere nimis prolixum nobis foret. Tandem postquam ad $\Sigma x^{p^\pi}$ perventum est, rei summa perfecta est. Namque posito
\[
P\equiv x^m+ax^{m-1}+bx^{m-2}+\text{etc.}
\]
erit $-a\equiv\Sigma x^{p^\pi}$, eodem modo coëfficiens secundus reliquarum functionum $(P,\rho^k)$, habebitur, unde reliqui ipsius $P$ determinari possunt. Saepius evenire potest, ut ad congruentias identicas perveniatur, ex quibus nihil derivari posse videtur. Quomodo huic difficultati obveniri possit, infra monstrabitur.

% A544-P0242-D03
\Article{364}
Omnia haec per exemplum multo clariora fient. Resolvenda proponitur functio $x^{15}-1$ secundum modulum $17$ in factores. Erit $m=4$ et quoniam productum ex omnibus functionibus elementaribus
\[
\equiv\frac{x^{15}-1\pdot x-1}{x^3-1\pdot x^5-1}=x^8-x^7+x^5-x^4+x^3-x+1
\]

\SourceLeaf{243}%
% A544-P0243-D01
Quare duo tantummodo erunt factores primi quatuor dimensionum $P$ et $P'$. Iam $x,xx,x^4,x^7,x^8,x^{11},x^{13},x^{14}$ in has duas periodos distribuantur
\[
\begin{aligned}
\Sigma x^{17^\alpha}&\equiv x+xx+x^4+x^8,\\
\Sigma x^{7\pdot17^\alpha}&\equiv x^7+x^{11}+x^{13}+x^{14}
\end{aligned}
\]
Sit secundum alteram functionem $P,P'$
\[
\Sigma x^{17^\alpha}\equiv A,\quad\Sigma x^{7\pdot17^\alpha}\equiv A'
\]
eritque
\[
\begin{aligned}
A+A'&\equiv1\\
AA&\equiv\Sigma x^{2\pdot17^\alpha}+\Sigma x^{3\pdot17^\alpha}+\Sigma x^{5\pdot17^\alpha}+\Sigma x^{9\pdot17^\alpha}\\
A'A'&\equiv\Sigma x^{14\pdot17^\alpha}+\Sigma x^{6\pdot17^\alpha}+\Sigma x^{5\pdot17^\alpha}+\Sigma x^{3\pdot17^\alpha}
\end{aligned}
\]
quare
\[
\begin{aligned}
AA+A'A'&\equiv\Sigma x^{17^\alpha}+\Sigma x^{7\pdot17^\alpha}+4\Sigma x^{3\pdot17^\alpha}+2\Sigma x^{5\pdot17^\alpha}\\
&\equiv1-4-4\equiv-7
\end{aligned}
\]
Hinc $A$ et $A'$ erunt radices congruentiae
\[
xx-x+4\equiv0\;(\mathrm{mod.}\,17)
\]
quae sunt $6,12$. Hinc $P$ dividet
\[
x^8+x^4+xx+x-6
\]
eritque
\[
\equiv x^4-6x^3-2xx-12x+1
\]
$P'$ autem erit $\equiv(P,\rho^7)$ eritque
\[
\equiv x^4-12x^3-2xx-6x+1
\]

% A544-P0243-D02
\Article{365}
Sufficit nobis hic possibilitatem solutionum harum monstravisse. Multa artificia, quibus hae operationes sublevari possunt, praeterimus brevitatis gratia. At consequentias quasdam pergraves praetermittere non possumus.

% A544-P0243-D03
Per praecedentia demonstratum est, omnes aequationes auxiliares pro solutione aequationis $x^\nu=1$, si in congruentias convertantur, habere radices possibiles, quando periodus
\SourceLeaf{244}%
% A544-P0244-D01
\[
x+x^p+x^{pp}+\dotstwo+x^{p^{m-1}}
\]
nondum est disiuncta. Subsistamus in casu, ubi $\nu$ est numerus primus; erit $m$ divisor ipsius $\nu-1$. Hic itaque congruentiae auxiliares, si numerus periodorum, quae per illas inveniuntur, est pars aliquota numeri $\frac{\nu-1}m$, habent radices reales. Si itaque $\frac{\nu-1}m$ est par i. e. si $m$ est divisor numeri $\frac{\nu-1}2$ seu si $p^{\frac{\nu-1}2}\equiv1\;(\mathrm{mod.}\,\nu)$ seu si $p$ est residuum quadraticum numeri primi $\nu$, aequatio quadratica, per quam radices in duas periodos dividuntur, habebit radices reales secundum modulum $p$. At in Cap. \textsc{vi} monstravimus, hanc aequationem posito $\nu=4n\pm1$ semper esse $xx+x\mp n=0$. Quare habetur insigne

% A544-P0244-D02
\textsc{Theorema.} \emph{Si numerus primus $p$ est residuum quadraticum numeri primi $4n\pm1$, congruentia}
\[
xx+x\mp n\equiv0(\mathrm{mod.}\,p)
\]
\emph{habebit radices reales, adeoque etiam congruentia}
\[
4xx+4x\mp4n\equiv0\quad\text{seu}\quad(2x+1)^2\mp\nu\equiv0
\]
\emph{i. e. $\pm\nu$ erit residuum quadraticum numeri $p$.}

% A544-P0244-D03
\Article{366}
Haec igitur est tertia theorematis fundamentalis Capitis \textsc{iv} completa demonstratio, eo magis attentione digna, quod principia, e quibus est petita, ab iis quibus ad priores usi sumus, prorsus sunt diversa.\ .At ex eodem hoc fonte, sed via opposita quartam deducamus. Scilicet sit $\nu$ numerus primus formae $4n\pm1$, $p$ alius primus quicunque, sitque $\pm\nu$ residuum quadraticum numeri primi $p$, demonstrabimus, $p$ fore residuum quadraticum numeri $\nu$.

% A544-P0244-D04
Sit $p^m$ minima potestas numeri $p$, quae sit $\equiv1(\mathrm{mod.}\,\nu)$. Divisores elementares functionis $\frac{x^\nu-1}{x-1}$ secundum $p$ habebunt $m$ dimensiones, quare omnium numerus erit $=\frac{\nu-1}m$. Iam quoniam $\pm\nu Rp$, congruentia
\[
xx+x\mp n\equiv0(\mathrm{mod.}\,p)
\]
erit resolubilis; sint radices $A,A'$. Distribuantur functiones $x,xx,\dotstwo x^{\nu-1}$ in binas classes per $\xi,\xi'$ designandas, erit
\SourceLeaf{245}%
% A544-P0245-D01
\[
\begin{aligned}
\xi+\xi'&\equiv A+A'+(1+x+xx+\dotstwo+x^{\nu-1})\\
\xi\xi'&\equiv AA'+\lambda(1+x+xx+\dotstwo+x^{\nu-1})
\end{aligned}
\]
quare
\[
(z-\xi)(z-\xi')-(z-A)(z-A')
\]
per quemvis divisorem elementarem functionis $\frac{x^\nu-1}{x-1}$ erit divisibilis. Hinc autem quivis horum divisorum elementarium aut $\xi-A$ et $\xi'-A'$, aut $\xi-A'$ et $\xi'-A$ dividet. Hinc patet (quoniam $A$ non $\equiv A'$), si pro $x$ ponatur $x^p$, $\xi$ et $\xi'$ non immutari. Si enim $\xi$ in $\xi'$ et vice versa transiret, $\xi-A$ et $\xi-A'$ per eandem functionem primam dividerentur. Q. E. A. Hinc denique sequitur, $\frac{\nu-1}2$ per $m$ dividi seu $p^{\frac{\nu-1}2}-1$ per $\nu$. Quare $p$ erit residuum quadraticum ipsius $\nu$. Q. E. D.

% A544-P0245-D02
Facile autem est omnes theorematis fundamentalis casus ex utroque theoremate derivare.

% A544-P0245-D03
\Article{367}
Quamvis ad casum, ubi $\nu$ est numerus primus, hic nos restrinxerimus, tamen etiam, si $\nu$ sit compositus, theoremata analoga haud magno negotio determinari possunt, quod fusius exponere brevitatis gratia nunc non licet.

% A544-P0245-D04
Manifestum est, similes observationes etiam de maiori periodorum multitudine formari posse. Ita si $\frac{\nu-1}m$ per $3$ dividitur i. e. si $p$ est residuum cubicum numeri primi $\nu$, aequatio, per quam radices aequationis $x^\nu=1$ in tres periodos distribuuntur quamque in Cap. \textsc{vi} a priori determinandam docuimus, solubilis erit secundum modulum $p$ et vice versa. Ita ex. gr. congruentia $x^3+xx-2x-1\equiv0$ secundum modulum primum quemcunque, qui est formae $7n\pm1$, resolvi potest, si vero aliam formam habeat, non poterit.

% A544-P0245-D05
Non difficile nobis foret hoc Caput multis aliis observationibus locupletare, nisi limites, intra quos restringi oportet, vetarent. Iis qui ulterius progredi amant, haec principia viam saltem addigitare poterunt.

% A544-P0245-D06
\Article{368}
Congruentiam aliquam $S\equiv0$ radices seu generalius divisores \emph{aequales} habere dicimus, si per functionis alicuius potestatem dividi possit.

